\documentclass[10pt,a4paper]{article}
\usepackage[T1]{fontenc}
\usepackage[margin=20mm]{geometry}
\usepackage{lmodern,amsmath,amssymb,amsthm,mathtools,microtype,xurl}
\usepackage[hidelinks]{hyperref}
\newtheorem{theorem}{Theorem}
\newtheorem{corollary}[theorem]{Corollary}
\title{Joint E/M obstruction at small period index}
\author{The Clankers}
\date{17 September 2026 --- Turn 1 research preprint}
\begin{document}\maketitle
\begin{abstract}
For the complete centered divisor box of an Erd\H{o}s--Straus shell we retain
both channel targets, their exact collisions, and the arithmetic relation
$p=4a\pmod R$. Every joint failure with effective period quotient of order at
most eight has an actually saturated inactive factor block after removal of
at most two prime occurrences. All such quotient profiles are classified.
The existing prime-shift index-six result is attributed; the composite-modulus,
cut-budget and coefficient-fibre extensions are proved here. No universal ES
occupancy or historical-priority claim is made.
\end{abstract}

\paragraph{Original source.}
Fix a prime $p\equiv1\pmod4$, $p/4<a<p/2$, and $R=4a-p$.
Write $a=\prod q_i^{e_i}$ and retain
\[
 \mathcal B=\prod[-e_i,e_i]_{\mathbb Z},\quad
 v(\beta)=\prod q_i^{\beta_i}\pmod R,\quad
 u=\prod q_i^{e_i+\beta_i},\quad
 \mu(g)=\#v^{-1}(g),\quad W=\operatorname{supp}\mu.
\]
The inverse is $\beta_i=v_{q_i}(u)-e_i$. The canonical M target is $-1$;
the E target is $-p^{-1}$, with the recorded alternative $-p$ under
$\beta\mapsto-\beta$. Thus joint failure is
$W\cap\{-1,-p,-p^{-1}\}=\varnothing$. The alternative is not a third channel.
These are the inherited Type I/II coordinates \cite{ET}.

Put $\Gamma=\langle-1,2,q_i\rangle$ in $G=(\mathbb Z/R)^\times$ and
$K=\{g:gW=W\}$. Since $1\in W$, every period lies in $W$, so $K\le\Gamma$
and this is also the actual stabilizer in $G$. In
\[
 A=\Gamma/K,\quad d=|A|,\quad S=W/K,\quad
 \alpha=-1K,\quad z=2K,\quad\eta=pK,\quad g_i=q_iK,
\]
retain the arithmetic equation
\[
 \boxed{\eta=z^2\prod_i g_i^{e_i}.}                       \tag{1}
\]
The extra coefficient unit $2$ has not been supplied as an available prime.
The ambient index $[G:K]$ can exceed the effective index $d$.

\begin{theorem}[Every target collision and factor cut]
On joint failure, $\alpha$ is a nonidentity involution, $d$ is even,
$\eta\ne\alpha$, and the distinct forbidden classes
$F=\{\alpha,\alpha\eta,\alpha\eta^{-1}\}$ number one, two, or three
according as $\eta=1$, $\eta^2=1\ne\eta$, or $\eta^2\ne1$.
Call $i$ active if $g_i\ne1$. Then
\[
 \operatorname{ord}(g_i)>2e_i+1,\qquad
 \boxed{2\sum_{i\text{ active}}e_i\le d-1-|F|.}        \tag{2}
\]
For a nonempty active-index set $J$, let $C_J=\langle g_i:i\in J\rangle$,
$U_J=\prod_{i\notin J}\{g_i^{-e_i},\ldots,g_i^{e_i}\}$ and
$F_J=(FU_J^{-1})\cap C_J$. Then
\[
 \boxed{2\sum_{i\in J}e_i\le |C_J|-1-\max(1,|F_J|).}    \tag{3}
\]
\end{theorem}
\begin{proof}
$S$ has no nonidentity period, since lifting one would contradict the definition
of $K$. It is symmetric under inversion. The collision assertions follow by
cancelling $\alpha$; $\eta=\alpha$ would put $1$ in $F$.
A factor whose interval fills its cyclic subgroup would give a period, so its
order exceeds $2e_i+1$. Translate its interval and adjoin $2e_i$ copies of
$\{1,g_i\}$. Each step strictly grows: otherwise $g_i$ would stabilize that
partial support and every subsequent product. This gives $|S|\ge1+2\sum e_i$,
proving (2). The same argument in $C_J$ gives the corresponding lower bound
for its subproduct. That subproduct avoids $F_J$ and cannot equal $C_J$, which
would also create a period. This proves (3), including the empty $F_J$ case.
\end{proof}
The single-target growth framework is classical Kneser theory; the two-target
prime-shift form already appears in \cite{Bryan}. No asymptotic estimate is used.

\newpage
\begin{theorem}[Complete classification through effective order eight]
Let $B=\prod_{q_i\in K}q_i^{e_i}$. If $d\le8$ and the shell fails jointly,
then
\[
 \boxed{W(B)=K,\qquad\Omega(a/B)\le2.}                   \tag{4}
\]
For nonpure cases the complete table is
\[
\begin{array}{c|c|c|c}
 A&\text{active capacity}&\text{active directions}&\text{phase}\\\hline
 C_6&1&\{g\},\ |g|=6&\eta=g^{\pm1}\\
 C_6&2&\{g,g^{-1}\},\ |g|=6&\eta=1\\
 C_8&1&\{g\},\ |g|=4&\eta=1\\
 C_8&1&\{g\},\ |g|=8&\eta=g^{\pm1}\\
 C_8&2&\{g,g^{-1}\},\ |g|=8&\eta=1\\
 C_4\times C_2&1&\{g\},\ |g|=4&\alpha\notin\langle g\rangle,\ \eta=g^{\pm1}.
\end{array}                                                  \tag{5}
\]
One occurrence means one first-power prime. Two mean one square prime or two
distinct simple primes, with their individual signs retained. The support is
$\{1,g^{\pm1}\}$ or $\{1,g^{\pm1},g^{\pm2}\}$ accordingly.
The pure cases $W=K$ are exactly: $C_2$ with $\eta=1$;
$C_2^2$ with independent $\alpha,z$ and $\eta=1$;
$C_6$ with $|z|=3$ or $6$;
$C_8$ with $|z|=8$;
and $C_4\times C_2$ with $|z|=4$ and $\alpha\notin\langle z\rangle$.
In the pure cases $\eta=z^2$. All rows retain (1) and generation of $A$ by
$\alpha,z,g_i$. Conversely, these data and $W(B)=K$ imply joint failure and
make $K$ the actual stabilizer.
\end{theorem}
\begin{proof}
At order two there is no active element by (2). At order four, an active
generator occurs once and leaves only $\alpha=g^2$ unoccupied, forcing
$\eta=1$; (1) instead makes $\eta$ an odd power of $g$. Thus the case is pure,
and its quotient cannot be cyclic since then $z^2=\alpha$.
At order six, every active element has order six, the budget permits at most
two occurrences, and all are in the same inverse pair. One gives the three-point
interval and odd $\eta$; excluding $\eta=\alpha$ leaves $g^{\pm1}$. Two give
the five-point interval missing only $\alpha$, forcing $\eta=1$.

At order eight, $C_2^3$ has no active element and cannot be generated purely
by $\alpha,z$. In $C_4\times C_2$, active elements have order four and capacity
one. Two in the same cyclic subgroup fill it, creating a period. Two in the
different cyclic subgroups give seven support elements, requiring $\eta=1$;
their central product lies outside the square subgroup, contradicting (1).
Three occurrences repeat a cyclic subgroup. For one, $z^2g=g^{\pm1}$, and
$\alpha$ must be outside its cyclic subgroup to avoid an exterior hit.

In $C_8$, mixed orders four and eight give seven support elements and an odd
$\eta$, whereas failure requires $\eta=1$. Two order-eight directions not in
one inverse pair fill the group. Three aligned occurrences give seven elements
but odd $\eta$. Two order-four occurrences fill their subgroup. Testing the
remaining one- or two-occurrence intervals gives precisely (5). The pure cases
follow because $A$ is generated by $\alpha,z$, with $\eta=z^2$.

Each active interval in (5) has an extreme quotient value attained by exactly
one active exponent vector. Its fine support is a translate of $W(B)$ and is
an entire nonempty $K$-coset. Therefore $W(B)=K$. For the converse each listed
support is aperiodic and avoids the indicated targets; multiplying by the
actual $K$-saturated block proves both assertions.
\end{proof}

Thus a non-subgroup joint failure first has $d=6$. A non-subgroup failure with
exactly two distinct target classes has $d\ge12$: no row up to eight allows
it, and $C_{10}$ has only the forbidden involution $\eta=\alpha$.
For $R=q\equiv3\pmod4$ prime with $(p/q)=-1$, an effective index-six failure
reduces further to $a=Br$ with exactly one first-power prime $r$ outside $K$,
$W(B)=K$, and $pK=r^{\pm1}K$. The existing ambient-index-six theorem is
credited to \cite{Bryan}; full ambient and effective indices are not conflated.

\newpage
\paragraph{The quotient does not erase the original coefficients.}
Choose $\sigma:A\to\Gamma$ with $\sigma(1)=1$. Every original unit has the
unique coordinates $(c,k)$ with $g=\sigma(c)k$, and the complete product is
\[
 (c,k)(d,l)=(cd,kl\omega(c,d)),\qquad
 \omega(c,d)=\sigma(c)\sigma(d)\sigma(cd)^{-1}\in K.       \tag{6}
\]
Expanding the representatives proves the normalized cocycle identity.
The inverse reads $c=gK$, $k=g/\sigma(c)$. The fine count remains
$\nu(c,k)=\mu(\sigma(c)k)$; its quotient is $\bar\mu(c)=\sum_k\nu(c,k)$.
The quotient kernel has the full star-difference basis
$e_{\sigma(c)k}-e_{\sigma(c)}$, $k\ne1$. Quotient values plus nonreference
coefficients reconstruct the original vector. The same statement holds on
each actual exponent fibre. No factor of $|K|$ is divided out of a count.

Writing $T=\tau(B^2)$, one active simple prime gives quotient masses $(T,T,T)$;
one active square gives five copies of $T$; two distinct active simple primes
give $T(1,2,3,2,1)$. The fine distributions retain the corresponding translates
of $\mu_B$, not uniform masses. A supported zero in this quotient retains its
source label and relation primitive, rather than becoming the absent $\tau$
of the original SplitZero scalar.

\paragraph{Actual arithmetic examples.}
For $p=3361$, $a=848=2^4\cdot53$, $R=31$,
\[
 K=\{1,2,4,8,16\},\quad d=6,\quad pK=53K.
\]
Both channels fail in this shell. The inactive coefficients are one at $1$
and two at the other four members of $K$. Each occupied quotient coset has
mass nine; dividing by five would not recover its integer coefficients.

For $p=5569$, $a=1397=11\cdot127$, $R=19$, the other phase occurs:
$K=\{1,7,11\}$ and $pK=127^{-1}K$. With $\sigma(j)=2^j$, $0\le j<6$,
\[
 \boxed{\omega(i,j)=7^{\lfloor(i+j)/6\rfloor}\pmod{19}.} \tag{7}
\]
Every lift of the quotient generator has order eighteen. Replacing the
extension by $C_6\times C_3$ would lose this actual cubic carry.

At $p=22129$, $a=5537=7^2\cdot113$, $R=19$, the exterior-only index-three
defect is rescued by M:
\[
 (u,h,r,s,\lambda)=(49,49,1,113,6),\qquad
 \frac4{22129}=\frac1{5537}+\frac1{735169638}+\frac1{6505926}.
\]
The composite shell $p=109,a=31,R=15$ has targets $(14,11,11)$ and neither
channel hits. The prime-parameter example $p=6340273,a=1259^2,R=51$ has
$W=K=\{1,35\}$, effective index eight and ambient index sixteen, demonstrating
why the two indices must remain separate.

\paragraph{Ordered reconstruction and scope.}
For any original hit, $d_0=\gcd(a,u)$ and
$(h,r,s)=(d_0^2/u,u/d_0,a/d_0)$ give $(a,hs\kappa,phr\kappa)$ in E with
$\kappa=(pr+s)/R$, or $(a,phs\lambda,phr\lambda)$ in M with
$\lambda=(r+s)/R$. The inverse is $u=a^2/(Ry-pa)$ or $pa^2/(Ry-pa)$,
respectively. The full workbench proves first-half completeness, every
normal-form case, all target cuts, and every original inverse. The standard-library
verifier and a separately implemented direct-divisor checker supply bounded
finite certificates; they do not prove universal occupancy or formal Lean status.

\begingroup\small
\begin{thebibliography}{3}
\bibitem{ET} Elsholtz, C., \& Tao, T. (2013). Counting the number of solutions to
the Erd\H{o}s--Straus equation on unit fractions. \emph{J. Aust. Math. Soc.,94}(1),
50--105. arXiv:1107.1010, \S2.
\bibitem{Bryan} chasebryan (2026, August 15). \emph{Two-target Kneser collapse
for external nonresidue shifts}, \S\S4--9; \emph{Exact index-six normal form
for the combined FAB targets}, \S\S3--8. \texttt{chasebryan/centl}, source blobs
\texttt{41e75319\ldots917d5}, \texttt{81f984e6\ldots{}b59e}. Full identifiers accompany
this source. The underlying addition theorem is DeVos (2013), arXiv:1303.3539,
Theorem 1, a proof of classical Kneser theory.
\end{thebibliography}
\endgroup
\end{document}
